\documentclass[11pt]{article}
\usepackage[margin=1in]{geometry}
\usepackage{amsmath,amssymb}
\usepackage{graphicx}
\usepackage{booktabs}
\usepackage{hyperref}
\usepackage{natbib}

\title{Caspase-11 Catalytic Activity and Autoprocessing Are Required for Noncanonical Inflammasome Speck Formation}
\author{Anonymous Authors}
\date{}

\begin{document}
\maketitle

\begin{abstract}
Caspase-11 (Casp11) directly binds cytosolic lipopolysaccharide (LPS) through its CARD domain and oligomerizes into a high-molecular-weight noncanonical inflammasome, but the relationship between its enzymatic activity and higher-order assembly has not been directly visualized. Here we develop a Casp11-mCherry fluorescent reporter and demonstrate that LPS transfection induces dose-dependent perinuclear speck formation in Casp11$^{-/-}$ bone marrow-derived macrophages (BMDMs) reconstituted with wild-type Casp11. The catalytically dead C254A mutant completely fails to form specks despite comparable expression, proving that CARD--LPS binding alone is insufficient for oligomeric assembly. In HEK293T cells, both catalytic activity (C254) and interdomain linker (IDL) cleavage at D285 are independently required for speck formation. Using an engineered TEV protease-cleavable Casp11 construct, we show that exogenous protease cleavage at the IDL rescues speck formation even in catalytically dead mutants, demonstrating that physical cleavage---not ongoing catalytic activity---is the critical event enabling higher-order oligomerization. Co-transfection experiments with untagged wild-type Casp11 and tagged C254A-mCherry show that wild-type rescues C254A speck formation through trans-cleavage, but this rescue requires cis-autoprocessing at D285 on the tagged molecule (the C254A/D285A double mutant is not rescued). We propose a revised model: LPS binding triggers Casp11 dimerization, dimerization activates catalytic activity, catalysis mediates intra-molecular autoprocessing at D285, and the processed form assembles into the higher-order oligomeric supramolecular organizing centre (SMOC).
\end{abstract}

\section{Introduction}

The noncanonical inflammasome pathway is activated when caspase-11 (murine) or caspase-4/5 (human) directly binds intracellular LPS from Gram-negative bacteria through their CARD domains \citep{shi2014inflammatory, man2017caspase4}. This binding triggers caspase oligomerization and activation, leading to cleavage of gasdermin D (GSDMD) and pyroptotic cell death \citep{kayagaki2015caspase11, shi2015cleavage}. The oligomeric caspase-11 complex is classified as a supramolecular organizing centre (SMOC) \citep{lee2018smoc}, analogous to the ASC speck of canonical inflammasomes.

Despite extensive study of caspase-11 signaling outputs, several fundamental questions about the assembly process remain unanswered. No one has directly visualized Casp11 SMOC formation in living macrophages. The field assumes that CARD--LPS binding alone drives assembly, but previous work showed that both catalytic activity (at C254) and IDL cleavage (at D285) are needed for GSDMD processing \citep{ross2018dimerization}. Whether these steps also control the upstream assembly of the inflammasome complex itself is unknown. Furthermore, it remains unresolved whether Casp11 autoprocessing occurs in cis or trans.

\section{Results}

\subsection{LPS Induces Dose-Dependent Casp11 Speck Formation in BMDMs}

Casp11$^{-/-}$ BMDMs transduced with Casp11 WT-mCherry via lentivirus formed large perinuclear mCherry specks upon LPS transfection. Speck formation was dose-dependent, increasing from $\sim$8\% of mCherry$^+$ cells at 0.5 $\mu$g/mL LPS to $\sim$35\% at 5 $\mu$g/mL. WT-mCherry also restored dose-dependent pyroptosis (LDH release up to $\sim$65\%) and GSDMD cleavage. The C254A catalytically dead mutant showed no speck formation, no cytotoxicity, and no GSDMD cleavage at any LPS dose despite comparable expression levels.

\subsection{Catalytic Activity Is Essential for SMOC Assembly}

In HEK293T cells, WT-mCherry formed specks spontaneously upon overexpression ($\sim$40--50\% of mCherry$^+$ cells) while C254A remained entirely diffuse ($\sim$2--5\%). Western blot showed that WT-mCherry was processed (cleaved bands visible) while C254A remained full-length. Blue native PAGE revealed that WT assembled into complexes exceeding 1~MDa, while C254A remained as monomer/dimer. This proves that CARD--LPS binding alone is insufficient; catalytic activity is essential for higher-order oligomerization.

\subsection{IDL Cleavage at D285 Is Independently Required}

The non-cleavable D285A mutant phenocopied the catalytic-dead C254A in failing to form specks ($\sim$2--5\%). The C254A/D285A double mutant also failed. This demonstrates that autoprocessing at the IDL is a separately required step for SMOC assembly, not merely a byproduct of caspase activity.

\subsection{TEV Protease Cleavage Rescues Speck Formation}

To distinguish whether catalytic activity is continuously required or whether physical cleavage is the activating event, we engineered a TEV cleavage site into the IDL of the D285A mutant. Co-expression with TEV protease rescued speck formation ($\sim$30--40\%) for TEV-D285A-mCherry. Critically, even the catalytically dead TEV-C254A/D285A-mCherry was rescued by TEV ($\sim$25--35\%), proving that physical cleavage at the IDL---not ongoing catalytic activity---is sufficient for oligomerization.

\subsection{Trans-Cleavage Followed by Cis-Dependent Oligomerization}

Co-transfection of increasing doses of untagged WT Casp11 with fixed amounts of tagged C254A-mCherry rescued C254A speck formation in a dose-dependent manner (from $\sim$2\% to $\sim$30\%). However, the C254A/D285A double mutant was not rescued, demonstrating that rescue requires cleavage at D285 on the tagged molecule itself.

Further dissection showed that untagged D285A (catalytically active but self-non-cleavable) also rescued C254A-mCherry processing and specks, while untagged C254A could not. This proves that the catalytic activity of the untagged partner cleaves C254A-mCherry in trans at D285, and the resulting processed molecule then assembles through cis-dependent conformational changes.

\section{Discussion}

Our findings establish a revised model for noncanonical inflammasome assembly. The process is not a single step triggered by LPS binding, but a multi-step pathway with distinct molecular requirements at each stage.

The key insight from the TEV rescue experiments is that physical cleavage at the IDL, delivered by any protease, is sufficient to enable higher-order oligomerization. This means that autoprocessing serves as a molecular switch: the conformational change induced by IDL cleavage exposes surfaces required for SMOC assembly. The intact IDL sterically or conformationally constrains the caspase, preventing the higher-order interactions needed for speck formation.

The trans-cleavage/cis-assembly paradigm has important implications. In a physiological context, initial LPS-induced dimers with catalytic activity can trans-cleave additional Casp11 molecules, which then undergo cis-dependent conformational changes to join the growing SMOC. This creates a positive feedback loop that amplifies the initial signal while maintaining the requirement for processing at each molecule.

\section{Methods}

Casp11-mCherry reporters were expressed via lentiviral transduction in Casp11$^{-/-}$ BMDMs or transient transfection (PEI) in HEK293T cells. Key mutants included C254A (catalytic dead), D285A (non-cleavable IDL), C254A/D285A (double mutant), and TEV-site variants. Cells were primed with Pam3CSK4 (4~h) and transfected with LPS from \textit{S.\ enterica} serovar Minnesota. Speck formation was quantified by confocal microscopy as percentage of mCherry$^+$ cells with at least one perinuclear speck. GSDMD cleavage and Casp11 processing were monitored by western blot. Cytotoxicity was measured by LDH release (16~h). Blue native PAGE assessed complex formation.

\bibliographystyle{plainnat}
\bibliography{references}

\end{document}
